% Principal GDZ PPN236018647,native488lower–489/printed481lower–482;HKUST489–490.
\begingroup\small
\begin{center}\bfseries BEMERKUNG.\end{center}
\indent Daß \textsc{Gauss} sich mit »Mannigfaltigkeiten von mehr als zwei Dimensionen« beschäftigt hat, zeigt eine Andeutung am Schluss der Selbstanzeige der \textit{Theoria residuorum biquadraticorum, commentatio secunda} vom 23. April 1831, Werke II, S. 178, und bestätigt ein Bericht von \textsc{Sartorius v. Waltershausen} über ein Gespräch mit \textsc{Gauss}, das zwischen 1847 und 1855 stattgefunden hat. »Wir können uns, sagte er, etwa in Wesen hineindenken, die sich nur zweier Dimensionen bewußt sind; höher über uns stehende würden vielleicht in ähnlicher Weise auf uns herabblicken, und er habe, fuhr er scherzend fort, gewisse Probleme hier zur Seite gelegt, die er in einem höhern Zustand später geometrisch zu behandeln gedächte« (\textsc{Gauss} zum Gedächtnis, Leipzig 1856, S. 81). Daß diese Gedanken auf eine frühe Zeit zurückgehen, wird durch einen Brief \textsc{Wachters} an \textsc{Gauss} vom 12. Dezember 1816 wahrscheinlich gemacht, der eine zwischen beiden Männern im April 1816 stattgehabte Unterhaltung widerspiegelt. \textsc{Wachter} schreibt\footnote{Siehe P. \textsc{Stäckel}, \textit{F. L. Wachter}, Mathematische Annalen \textbf{54}, 1901, S. 49, wo S. 61—69 der ganze Brief veröffentlicht ist. Werke VIII, S. 175 sind nur die auf die anti-Euklidische Geometrie bezüglichen Stellen abgedruckt.}:

\indent »Vor kurzem suchte ich die Construction des Fundamentalkörpers, (das, was der Würfel für den Raum von 3 Dimensionen) allgemein für einen Raum von beliebig vielen Dimensionen. Ich fand, daß dieses, \textls[120]{wofern eine gewisse Schlußart richtig ist}, mit der Zahl der rechten Winkel um einen Punkt in der Ebene zusammenhänge, daß nämlich der Raum \(n\) Dimensionen haben würde, wenn die Summe der rechten Winkel um einen Punkt herum \(=2^{n-1}\), und daß dann die Zahl der \textls[120]{Ecken} des Fundamentalkörpers in diesem Raum von \(n\) Dimens[ionen] durch \(2^n\), die Zahl der Begränzungen der ersten Art von der Dimension \((n-1)\) für jede Ecke durch \(n\), und die Zahl derselben, der Gränzkörper von der Dim[ension] \((n-1)\), für den ganzen Körper durch \(2n\), die Zahl der Begränzungen\begingroup\makeatletter\def\@makefnmark{\hbox{\normalfont\@thefnmark}}\makeatother\renewcommand{\thefootnote}{[**)]}\footnotemark\renewcommand{\thefootnote}{[**)}\footnotetext{In der Handschrift steht \textls[120]{Berührungen}.]}\endgroup\ der zweiten Art von der Dim[ension] \((n-2)\) durch \(4n\), sein körperlicher Inhalt, wenn die Seite \(L\), durch \(L^n\), und seine Diagonale durch \(L\sqrt{n}\) gegeben werde. —
\par\vfill\noindent X1.\hfill61
\clearpage\SourceLeaf{489}\ArithmeticHead{482}{L}{GEOMETRIE. NACHLASS.}
\noindent Von einer andern Seite läßt sich zeigen: daß \textls[120]{wenn es möglich, daß eine gerade Linie Zweige habe}, sie derselben unendlich viele, und \textls[120]{auch der Raum unendlich viele Dimensionen} haben würde. Auf \textls[120]{gewisse Weise} könnte man sagen, daß auch in Ihrer, der anti-euklidischen Geometrie, der Raum unendlich viele Dimensionen habe, die aber alle wieder im unendlichen liegen. Nämlich es sei \(BAC\) ein rechter Winkel, \(AD\) die Constante für das rechtwinklichte asymptotische Dreieck, dessen anderer Winkel \(=45^0\); der rechte Winkel \(BAC\) werde durch \(AD\) halbiert, dann ließe sich dieselbe Construction durch \(DE,EF\) u.s.w. an den Punkten \(D,E,F\) u.s.w. wiederholen, ohne daß die Linie \(AC\) jemals erreicht würde, — und dieses auf analoge Art räumlich gedacht gäbe innerhalb des von drei rechtwinklichten Ebenen begränzten Raumes körperliche Ecken, welche zuerst durch 3 Ebenen von \(45^0\) Oeffnung begränzt, der Reihe nach durch eben so viel Körper von der Dimension 3, 4, 5 u.s.w. bis \(\infty\), begränzt gedacht werden könnten.«

\indent Auch der Brief von \textsc{Gauss} an \textsc{Grassmann} vom 14. Dezember 1844 (siehe oben S. 436) ist in diesem Zusammenhang zu erwähnen; sagt doch \textsc{Gauss} darin, dass \textsc{Grassmanns} Tendenzen teilweise denjenigen Wegen begegneten, auf denen er selbst nun seit fast einem halben Jahrhundert gewandelt sei, und bezeichnet er als ein Zeugnis für diese Bestrebungen die oben erwähnte Selbstanzeige vom Jahre 1831. Man vgl. hierzu auch noch die Andeutungen im art. 5. der \textit{Beiträge zur Theorie der algebraischen Gleichungen} vom Jahre 1849, Werke III, S. 79\footnote{Die von \textsc{Gauss} begehrte Lehre von den »nach der Stetigkeit zusammenhängenden Größencombinationen« hat \textsc{Riemann} in dem Habilitationsvortrag vom 10. Juni 1854 begründet, \textsc{Riemanns} Werke, 2. Aufl., S. 272. Bei der Entwicklung des Begriffs mehrfach ausgedehnter Größen verweist \textsc{Riemann} (a. a. O. S. 273 ausdrücklich auf die oben im Text genannten Veröffentlichungen von \textsc{Gauss}.}.

\indent Eine weitere, ausgiebigere Quelle ist A. \textsc{Ritter} (1826—1908), der von Michaelis 1850 bis 1853 in Göttingen studiert hat und 1853 unter \textsc{Gauss} mit einer Abhandlung \textit{Ueber das Princip des kleinsten Zwanges} promoviert worden ist. Daß die unter [I.] abgedruckten, den siebenten Abschnitt der Dissertation bildenden Ausführungen auf Anregungen von \textsc{Gauss} zurückgehen, würde nur den Wert einer Vermutung haben, wenn nicht der auf der Bibliothek der Technischen Hochschule zu Aachen aufbewahrte handschriftliche Nachlaß \textsc{Ritters} die Ausarbeitung einer Vorlesung über die Methode der kleinsten Quadrate enthielte, die dieser im Wintersemester 1850/51 bei \textsc{Gauss} gehört hat. Den Schluß der Handschrift bildet nämlich die unter [II.] abgedruckte Notiz, die, wie es scheint, den Inhalt der beiden letzten Vortragstunden wiedergibt.

\indent Ob \textsc{Ritter} die Gedanken von \textsc{Gauss} überall richtig aufgefasst hat, muß dahingestellt bleiben; immerhin macht gerade die Stelle, die sich auf Mannigfaltigkeiten von \(n\) Dimensionen bezieht, einen Vertrauen erweckenden Eindruck. Die Behandlung des Minimumsproblems ist in verschiedener Beziehung unvollständig. Es genüge hier zu bemerken, daß man nicht immer Punkte finden kann, für die \(n\) der Ausdrücke \(u_1,.\ .\ .,u_m\) verschwinden, während gleichzeitig die andern \(m-n\) positiv ausfallen. Zum Beispiel verschwinden für \(n=3\) und \(m=8\), wenn es sich um das Innere eines von acht Dreiecken begrenzten Achtflächers handelt, in jeder der sechs Ecken vier der Ausdrücke \(u_1,.\ .\ .,u_8\). Im übrigen vergleiche man den in der zweiten Abteilung dieses Bandes abgedruckten Aufsatz »\textsc{Gauss} als Geometer«.
\par\noindent\hfill\textsc{Stäckel}.
\ClosingRule
\endgroup
